\documentclass[12pt]{article}
\usepackage[left=0.5in, right=0.5in, top=1in, bottom=1in]{geometry}
\usepackage{authblk}


\title{Transition Costs and Increasing Rotational Complexity}
\author{}
%\author{Lawson Connor}
%\author{Victor Funes-Leal}
%\author{Eunchun Park}
%\affil{Department of Agricultural Economics and Agribusiness\\ University of Arkansas}

\date{}
\begin{document}

\maketitle

\begin{abstract}

There is growing interest in promoting climate-smart practices, such as more complex crop rotation. Much of this advocacy is based on evidence that more diverse rotations can improve soil health and raise average yields. What is often missing from these discussions is the opportunity cost of transitioning away from monoculture, as well as its implications for farm profitability and risk. In practice, adoption decisions are not binary. They involve changes in crop mix, marketing channels, and exposure to both yield and price variability over time.

We utilize field-level data from Illinois on crop histories and yields, combined with grain-elevator-specific cash bids, to estimate how six-year rotation sequences affect expected yields, yield variability, and net economic returns for corn and soybeans. Using local cash prices allows us to account for spatial differences in market access, such as proximity to ethanol plants and other demand centers. With these data, we estimate the net present value (NPV) of alternative transition paths from monoculture to more complex rotations, explicitly accounting for foregone revenue from planting a different crop than the focal crop.

Our data indicate that approximately 9 percent of the sample’s acreage is farmed in monoculture, a pattern that appears to be related, in part, to market access and price incentives. Our results suggest that transitioning away from monoculture toward perfect rotation increases yields at each point along the transition path. However, these gains are accompanied by increased yield variability, which persists even after dispersion is normalized using the coefficient of variation. Moreover, the yield gains from rotation must be weighed against the opportunity cost of replacing corn with soybeans (or vice versa), which can materially affect the profitability of transition paths. As a result, the rotations that maximize expected yield are not necessarily those that maximize net economic returns.

Using local elevator cash bids, we aim to determine whether yield gains from increased rotation are sufficient to justify transition costs after accounting for changes in yield risk and marketing opportunity costs. While rotation improves average productivity, it also raises yield dispersion and can lower net returns once foregone revenue from alternative crop choices is considered. This complicates the common framing of rotation as a low-cost resilience investment.

The paper concludes by discussing how these results bear on conservation incentives, crop insurance design, and the practical barriers to adoption. In particular, we highlight the challenges of encouraging rotation through policies that implicitly treat adoption as a binary decision and ignore transition costs, nonlinear yield responses, and local market conditions. Limitations of the current project include the use of total variance rather than downside risk and the omission of changes in input costs that may result from transitioning. These results are intended to generate discussion about how climate-smart practice incentives can be better aligned with the economic realities farmers face when making long-run production decisions about practice adoption.

\end{abstract}

\end{document}